The derivation of \(G\) takes that reciprocity to a much more concrete geometric condition.\par

Dimensional Mechanics first determines what type of object \(G\) can be. It fixes its magnitude line and dimensional homogeneity. Chronological closure then selects its value: dimension determines the class of magnitude, and closure determines its dimensionless coefficient.\par

The selection proceeds from the temporal periodicity of (108) phases.\par

The cycle \(108\) determines a frequency, a phase radius, and an energy. From that energy arises a chronological mass. And constructing the reduced Compton wavelength of that mass yields exactly the same radius as that of the cycle:\par

\[
\bar\lambda_C=R_{108}.
\]

Phase and matter are already stating the same length before \(G\) is introduced.\par

Then gravity comes into play. The reduced gravitational radius is constructed.\par

\[
r_g=\frac{Gm}{c^2},
\]

and gravitational closure is required to identify the phase radius, Compton
length, and gravitational radius as a single structural scale:\par

\[
r_g=R_{108}=\bar\lambda_C.
\]

That condition uniquely selects\par

\[
\theta_{108}
=
\left(\frac{54}{\pi}\right)^2.
\]

In this way,\par

\[
R_{108}
=
\bar\lambda_C
=
r_g
=
\ell_P.
\]

The structural function of \(\theta_{108}\) determines the significance of its terminal decimal: it shows what \(G\) is doing.\par

\(G\) appears as the coefficient required to make the four lengths published by these languages coincide:\par

\begin{itemize}[leftmargin=*,itemsep=0.35em]
\item the lenguaje of the fase produce \(R_{108}\);
\item the quantum language produces \(\bar\lambda_C\);
\item the gravitational language produces \(r_g\);
\item the lenguaje of Planck produce \(\ell_P\).
\end{itemize}

The gravitational constant is the mediator that prevents those four readers from separating.\par

Stated in more physical terms: \(G\) expresses how strongly geometry must respond for a quantum whose mass follows from the clock \(108\) to close upon its own phase.\par

It is the relation that makes possible for matter, phase, and geometry to be three compatible publications of the same event.\par

Here the Machian reading becomes stronger. The local gravitational scale is selected by a global cycle-closure condition. The local radius of a quantum acquires meaning when integrated into a complete periodicity. Inertia, phase, and geometry are determined relationally.\par

The demonstrated scope provides a concrete mathematical mechanism for a Machian reading: a local property—the gravitational coupling—is selected by the consistency of the global state. The other historical formulations of Mach's principle preserve their own domain.\par

Then comes the fifth dimension.\par

The scalar \(G\) fixes the gravitational intensity, and the carrier specifies its representational structure. HMT constructs a symmetric traceless tensor space of dimension five. That carrier admits five helicoidal weights:\par

\[
+2,\quad +1,\quad 0,\quad -1,\quad -2.
\]

They are the five natural degrees of freedom of a symmetric traceless tensor in three dimensions before the dynamical constraints of a specific theory are imposed.\par

That clarifies a point that is often confused. “Five components” designates the complete carrier, whereas the two helicities designate the transverse massless reduction in general relativity.\par

It means that the complete carrier has five irreducible coordinates. The dynamics then decides which part of that carrier is physical:\par

\begin{itemize}[leftmargin=*,itemsep=0.35em]
\item in a massive spin-two realization, all five weights may survive;
\item in a massless realization with gauge invariance, the transverse traceless projection eliminates the sectors \(\pm1\) and \(0\);
\item the helicities \(\pm2\) remain.
\end{itemize}

Thus, HMT provides a space prior to the dynamical decision. It constructs the complete five-component carrier and contains, through the corresponding projection, its physical reduction to two helicities.\par

Beauty resides in the ordered coexistence of two statements:\par

\begin{itemize}[leftmargin=*,itemsep=0.35em]
\item the spin-two geometry has a five-component carrier;
\item massless gravitational radiation occupies a two-dimensional subspace.
\end{itemize}

Both statements form a coherent and precise reduction.\par

And that reduction preserves another central lesson: the scalar \(G\) fixes the common intensity, while the orientation information resides in the tensor carrier. The intensity resides in the constant, the helicities reside in the representation, and the complete genealogy brings magnitude and carrier together.\par

